\chapter{Implementation}\label{ch09:implementation}

The preceding chapters developed the mathematics of the typed finite-function
kernel: objects with ordered carriers, total arrows with declared types,
endomap monoids and their factorial idempotent powers, semigroupoid paths,
the deterministic standard-basis lift with its exact norm theory, and
shortlex denotational normal forms. This chapter documents the executable
artifact that realizes those constructions: the Python package
\texttt{src/tks\_v8\_core}, comprising three substantive modules ---
\texttt{category.py}, \texttt{deterministic\_lift.py}, and
\texttt{normal\_forms.py} --- together with a small \texttt{\_\_init\_\_.py}
that re-exports the public names.

The documentary stance of this chapter is strict: we describe the code as it
is at the manifested snapshot, not as an idealization of it. Every class,
constructor precondition, rejection behavior, and error type reported below
is read directly from the source files whose SHA-256 digests are bound in the
prototype manifest (Section~\ref{ch09:sec:manifest}). Where a design decision
is visible in the code --- a field excluded from equality, a guard placed
before rather than after an insertion, a name generated by hashing rather
than by decimal expansion --- we record the decision and the rationale that
the code and its docstrings themselves supply. The artifact carries the
status \textsc{gate 1 prototype / in review / not frozen}, admission Ring~4,
license status \textsc{not licensed}; nothing in this chapter promotes that
status.

\section{The package and its documentary boundaries}\label{ch09:sec:package}

\subsection{Module inventory}

The package consists of four files:

\begin{center}
\begin{tabular}{@{}llr@{}}
\toprule
File & Role (as in the manifest) & Bytes \\
\midrule
\texttt{src/tks\_v8\_core/\_\_init\_\_.py} & package exports & 1{,}018 \\
\texttt{src/tks\_v8\_core/category.py} & typed finite-function category & 9{,}850 \\
\texttt{src/tks\_v8\_core/deterministic\_lift.py} & deterministic standard-basis lift & 5{,}483 \\
\texttt{src/tks\_v8\_core/normal\_forms.py} & denotational shortlex closure & 9{,}889 \\
\bottomrule
\end{tabular}
\end{center}

The dependency structure is a chain: \texttt{deterministic\_lift.py} imports
\texttt{FiniteArrow} and \texttt{ValidationError} from \texttt{category.py};
\texttt{normal\_forms.py} imports \texttt{CompositionTypeError},
\texttt{FiniteArrow}, \texttt{FiniteObject}, \texttt{ValidationError}, the
module-internal helper \texttt{\_ordered\_tuple}, and the constructors
\texttt{compose} and \texttt{identity} from \texttt{category.py}. Neither of
the two downstream modules imports the other, and no module imports anything
outside the Python standard library (\texttt{collections.abc},
\texttt{dataclasses}, \texttt{hashlib}, \texttt{math}, \texttt{numbers},
\texttt{heapq}, \texttt{itertools}, \texttt{typing}). There is no dependency
on the production CAS notation compiler, on any numerical library, or on any
network or filesystem resource. This is not incidental hygiene: the
manifest's integration boundary asserts
\texttt{production\_runtime\_dependency: false}, and the boundary regression
test of Chapter~10 asserts the same field.

\subsection{What the modules deliberately do not contain}

\begin{boundary}\label{ch09:bdy:no-semantics}
The module docstring of \texttt{category.py} states the boundary in its
second sentence: the module ``does not assign CAS or Noetic meaning to an
arrow.'' No identifier anywhere in the three modules names a Noetic, a
Domain, a Drive, a Register, an RPM pathway, or any other construct of the
TKS/CAS ontology. The kernel supplies a reference semantic layer in which an
arrow always carries an exact source and target and composition is rejected
when those types do not match --- nothing more. The mapping of canonical
content onto this substrate is proof obligation PO-TKS8-CORE-08, which is
\textsc{open / not licensed} (Chapter~12).
\end{boundary}

\section{\texttt{category.py}: the typed finite-function category}\label{ch09:sec:category}

\subsection{Purpose}

\texttt{category.py} implements the category $\cC$ of Chapters~3 and~4 as
executable data: named finite objects, total finite-function arrows with
declared source and target, partial type-checked composition, the endomap
wrapper that alone carries monoid powers, and the semigroupoid path that
alone carries heterogeneous typed composites. The module realizes, in fewer
than ten kilobytes, the three disciplines that the red-team audit found
missing from kind-level checking: typed composition, power discipline, and
deterministic identity.

\subsection{Error taxonomy}\label{ch09:sec:cat-errors}

The module defines exactly two exception types, and their placement in
Python's built-in exception hierarchy is itself a design statement:

\begin{center}
\begin{tabular}{@{}lll@{}}
\toprule
Exception & Base class & Meaning \\
\midrule
\texttt{ValidationError} & \texttt{ValueError} & malformed object, arrow, or path \\
\texttt{CompositionTypeError} & \texttt{TypeError} & source/target objects do not match \\
\bottomrule
\end{tabular}
\end{center}

\texttt{ValidationError} subclasses \texttt{ValueError}: the caller supplied
data of an acceptable Python type whose \emph{value} violates a kernel
precondition (a blank name, a duplicate carrier element, an out-of-range
image index, a negative exponent). \texttt{CompositionTypeError} subclasses
\texttt{TypeError}: the caller attempted an operation that is ill-typed
\emph{categorically} --- composing arrows whose middle objects differ, or
wrapping a non-endo arrow as an \texttt{Endomap}. The distinction mirrors
the mathematical one between a malformed datum and a well-formed datum used
at the wrong type, and it lets test fixtures assert the precise failure
class rather than a generic exception. Both exception types carry
human-readable messages that name the offending arrow or object by its
presentation name and spell out both type signatures in the composition
case, so a rejection is diagnosable from the message alone.

\subsection{The ordered-input guard}\label{ch09:sec:ordered-tuple}

The private helper \texttt{\_ordered\_tuple(values, field)} is the single
chokepoint through which every ordered input in the package passes. Its
docstring reads: ``Reject containers whose iteration order is not part of
their contract.'' The implementation is a single guard:

\begin{quote}
\texttt{if (isinstance(values, (str, bytes, bytearray, Mapping, Set)) or not isinstance(values, Sequence)): raise ValidationError(...)}
\end{quote}

followed by \texttt{return tuple(values)}. The accepted inputs are exactly
the objects registered as \texttt{collections.abc.Sequence} --- lists,
tuples, and other sequence types with a stable, replayable order --- minus
the string-like sequences \texttt{str}, \texttt{bytes}, and
\texttt{bytearray}, which are technically sequences but almost always
signal a caller error when passed where a sequence of elements is expected.
Rejected are: Python \texttt{set} and \texttt{frozenset} values (whose
iteration order depends on insertion history and on
\texttt{PYTHONHASHSEED}), mappings (whose iteration yields keys, and whose
acceptance would silently discard values), and one-shot iterables such as
generators (which are not \texttt{Sequence} instances, and whose consumption
would make a second validation pass observe different data).

\begin{ruleenv}[Ordered-input rejection]\label{ch09:rule:ordered}
Carrier elements, arrow mappings, path arrow sequences, word label
sequences, catalog object sequences, catalog generator sequences, and lift
coordinate vectors are all validated through this guard (or, for the two
lift vector entry points, through an inline copy of the same check).
A set, mapping, or one-shot iterable at any of these positions raises
\texttt{ValidationError} at construction, before any downstream computation
can observe hash-order data.
\end{ruleenv}

The error message is deliberately uniform --- ``\texttt{<field> must be an
ordered sequence; sets, mappings, and one-shot iterables are not
accepted}'' --- with the field name interpolated, so every rejection site
reports both what was wrong and where. The rationale, developed at length in
Chapter~3, is compressed into the specification's one sentence: the order of
these inputs changes basis indices, typed denotations, or shortlex rank, so
silently accepting hash-order iteration would make results depend on
interpreter state. The guard converts a reproducibility hazard into an
immediate, attributable construction error.

\subsection{Compact integer labels}\label{ch09:sec:compact-label}

The second private helper, \texttt{\_compact\_integer\_label(value)},
renders a nonnegative integer for use inside a presentation name. For
integers of at most 256 bits it returns the ordinary decimal string. For
larger integers it encodes the value big-endian into bytes, takes the first
sixteen hexadecimal characters of the SHA-256 digest of that encoding, and
returns the string \texttt{bits=<bit\_length>,sha256=<digest16>}.

Two facts make this helper load-bearing rather than cosmetic. First, CPython
by default refuses to convert integers with more than a few thousand digits
to decimal strings (the integer-string conversion length limit), so naive
\texttt{str(exponent)} inside a power name would raise for very large
exponents; the helper's cost is $O(\log n)$ in the exponent and cannot hit
that limit. Second, the label appears only in presentation names, which
Definition~4.\emph{x} of the arrow discipline (and the field declaration
\texttt{compare=False} documented below) exclude from extensional equality
--- so the truncated digest cannot introduce a collision into any equality,
hashing, or denotation computation. The regression suite exercises exponents
$2^{30}$ (below the threshold, decimal path) and $10^{5000}$ (above it,
digest path).

\subsection{\texttt{FiniteObject}}\label{ch09:sec:finiteobject}

\begin{construction}[The object type]\label{ch09:con:object}
\texttt{FiniteObject} is a frozen dataclass with \texttt{slots}, with two
fields: \texttt{name: str} and \texttt{elements: tuple[Hashable, ...]}.
Construction runs a \texttt{\_\_post\_init\_\_} validator that enforces, in
order:
\begin{enumerate}
\item \texttt{name} is a \texttt{str} whose \texttt{strip()} is nonempty
      (``a finite object needs a non-blank name'');
\item \texttt{elements} passes the ordered-input guard of
      Rule~\ref{ch09:rule:ordered} and is frozen to a tuple;
\item every element is hashable --- checked by attempting to build a
      \texttt{set}; a \texttt{TypeError} from unhashable elements is caught
      and re-raised as \texttt{ValidationError} with the original exception
      chained as its cause;
\item the elements are pairwise distinct --- the set's cardinality must
      equal the tuple's length.
\end{enumerate}
\end{construction}

The public surface beyond the two fields is small: \texttt{index(element)}
returns the standard-basis index of an element, converting the underlying
\texttt{ValueError} of tuple search into a \texttt{ValidationError} whose
message names both the element and the object; and \texttt{\_\_len\_\_}
returns the carrier size, so $|X|$ is written \texttt{len(x)} in client
code.

\begin{invariant}[Name-inclusive object identity]\label{ch09:inv:object-id}
Because \texttt{FiniteObject} is a frozen dataclass with both fields
participating in equality and hashing, two objects are equal exactly when
their names and their carrier tuples (with order) are equal. The class
docstring states the consequence: two carriers with the same elements but
different declared object identities ``are not silently interchangeable.''
Every downstream type check --- composition legality, endomap
admissibility, path segment matching, catalog membership --- inherits this
identity, so the name participates in typing everywhere, uniformly, through
ordinary equality.
\end{invariant}

Freezing matters beyond immutability hygiene: frozen dataclasses are
hashable, and \texttt{normal\_forms.py} uses objects directly as dictionary
keys (in the object-rank table and inside denotation keys). The
\texttt{slots} declaration additionally prevents attribute injection ---
no caller can attach ad hoc metadata to an object and smuggle state past
the declared fields.

\subsection{\texttt{FiniteArrow}}\label{ch09:sec:finitearrow}

\begin{construction}[The arrow type]\label{ch09:con:arrow}
\texttt{FiniteArrow} is a frozen, slotted dataclass with four fields:
\texttt{name} (declared \texttt{field(compare=False)}), \texttt{source} and
\texttt{target} (both \texttt{FiniteObject}), and \texttt{mapping}, a tuple
of integers in which \texttt{mapping[i]} is the target-basis index of the
image of source element $i$. Validation enforces, in order:
\begin{enumerate}
\item a non-blank presentation name;
\item \texttt{source} and \texttt{target} are \texttt{FiniteObject}
      instances (``an arrow source and target must be FiniteObject
      values'');
\item \texttt{mapping} passes the ordered-input guard;
\item the mapping length equals the source carrier length --- totality; the
      message counts the required images explicitly (``needs one image for
      each of the $n$ source elements'');
\item every image index is an \texttt{int} that is \emph{not} a
      \texttt{bool} (the guard \texttt{isinstance(image\_index, bool)} is
      checked first, because Python's \texttt{bool} subclasses
      \texttt{int} and would otherwise pass silently as $0$ or $1$);
\item every image index lies in the half-open range $[0, |Y|)$; the message
      names the offending index and the target object.
\end{enumerate}
\end{construction}

\begin{invariant}[Extensional arrow equality]\label{ch09:inv:extensional}
Because \texttt{name} is declared with \texttt{compare=False}, dataclass
equality and hashing are computed from \texttt{(source, target, mapping)}
only. Two arrows with different presentation names but identical typed
mapping data are equal and hash identically; conversely, arrows differing
in source, target, or any image index are unequal even under the same name.
This is the executable form of the extensionality discipline of Chapter~4:
names are for humans and diagnostics, denotations are for mathematics.
\end{invariant}

The asymmetry with \texttt{FiniteObject} --- where the name \emph{does}
participate in equality --- is deliberate and worth pausing on. An object's
name is part of its type identity: two registers that both carry
$\{0,1\}$ are different objects, and conflating them would license exactly
the silent coercions the kernel exists to reject. An arrow's name, by
contrast, is a presentation artifact: the same function $f : X \to Y$ may be
constructed under many names (as a composite, as a power accumulator, as a
closure candidate), and all of those constructions must be recognized as the
same denotation for the normal-form machinery to work at all. The code
encodes the distinction in one keyword argument.

Three further members complete the arrow surface:
\begin{itemize}
\item \texttt{\_\_call\_\_(element)} applies the arrow to a carrier
      element: it resolves the element's source index via
      \texttt{source.index} (raising \texttt{ValidationError} for foreign
      elements) and returns the target carrier element at the mapped index.
      Arrows are thus directly usable as functions on carrier elements, not
      only on indices.
\item \texttt{is\_endomap} is a property testing \texttt{source == target}
      --- name-inclusive object equality, per
      Invariant~\ref{ch09:inv:object-id}. An arrow between two distinct
      same-carrier objects is \emph{not} an endomap.
\item \texttt{denotation\_key} is a property returning the triple
      \texttt{(source, target, mapping)} --- ``a name-independent key for
      the denoted finite function,'' in the docstring's words. This triple
      is precisely the equality-participating data, and
      \texttt{normal\_forms.py} uses it as the dictionary key for the
      closure table.
\end{itemize}

\subsection{Identity and composition}\label{ch09:sec:compose}

Two module-level constructors realize the categorical structure.

\texttt{identity(obj, *, name=None)} builds the identity arrow on an
object: source and target both \texttt{obj}, mapping
\texttt{tuple(range(len(obj)))}, default name \texttt{id[<object name>]}.

\texttt{compose(after, before, *, name=None)} returns the categorical
composite ``\texttt{after} after \texttt{before}.'' Its docstring is
explicit about the guard: it returns the composite ``after checking the
middle object exactly.'' The check is

\begin{quote}
\texttt{if before.target != after.source: raise CompositionTypeError(...)}
\end{quote}

where the inequality is full object inequality --- name and ordered carrier
--- and the error message prints both arrows with their complete type
signatures (\texttt{f:X->Y then g:Z->W}), making the mismatch
self-describing. When the check passes, the composite is a fresh
\texttt{FiniteArrow} with source \texttt{before.source}, target
\texttt{after.target}, mapping

\begin{quote}
\texttt{tuple(after.mapping[index] for index in before.mapping)}
\end{quote}

and default name \texttt{(<after> o <before>)}. The mapping expression is
the whole of composition: index chasing through two tuples, with no
arithmetic, no floating point, and no reference to carrier elements. The
argument order --- \texttt{after} first --- matches the mathematical
notation $g \circ f$ read aloud, and the keyword-only \texttt{name}
parameter keeps presentation naming an explicit, optional act.

\begin{remark}\label{ch09:rem:no-coercion}
Nothing in \texttt{compose} inspects carrier contents or cardinalities
except through object equality. There is no code path that coerces a
mismatch: no ``compatible carrier'' fallback, no implicit reindexing, no
name-blind comparison. PO-TKS8-CORE-02 (reject every source/target
mismatch) is discharged at the level of executable validation by this
single guard plus the negative tests of Chapter~10 that drive foreign
objects, same-carrier-different-name objects, and reversed-order arguments
into it.
\end{remark}

\subsection{\texttt{Endomap}: certified membership in one monoid}\label{ch09:sec:endomap}

\begin{construction}[The endomap wrapper]\label{ch09:con:endomap}
\texttt{Endomap} is a frozen, slotted dataclass wrapping a single field
\texttt{arrow: FiniteArrow}. Its validator enforces two conditions: the
wrapped value is a \texttt{FiniteArrow} (else \texttt{ValidationError}),
and \texttt{arrow.is\_endomap} holds (else \texttt{CompositionTypeError},
with a message that prints the arrow's actual type signature and states
``so it is not an endomap''). The docstring calls the wrapper ``an arrow
certified to lie in one endomorphism monoid End(X)'' --- construction is
the certificate.
\end{construction}

The property \texttt{carrier} returns the (shared) source object. The
substantive method is \texttt{power(exponent)}:

\begin{enumerate}
\item \emph{Exponent validation.} The exponent must be an \texttt{int},
      not a \texttt{bool}, and nonnegative; otherwise
      \texttt{ValidationError} (``an endomap exponent must be a
      non-negative integer''). Zero is legal and yields the identity.
\item \emph{Binary powering.} The method initializes an accumulator to
      \texttt{identity(carrier)} under the internal name
      \texttt{power-accumulator}, sets the base to the wrapped arrow, and
      loops over the bits of the exponent: when the low bit is set, the
      accumulator becomes \texttt{compose(base, accumulator)}; the exponent
      is shifted right; and --- only if bits remain --- the base is squared
      via \texttt{compose(base, base)} under the internal name
      \texttt{power-square}. Guarding the final squaring on
      \texttt{remaining} being nonzero avoids one wasted composition per
      call, a small courtesy visible in the code.
\item \emph{Result packaging.} The method returns a fresh
      \texttt{Endomap} whose arrow carries the accumulator's mapping and
      the presentation name
      \texttt{<base name>\^{}(<compact label>)}, with the exponent rendered
      by \texttt{\_compact\_integer\_label}
      (Section~\ref{ch09:sec:compact-label}).
\end{enumerate}

Because each loop iteration performs at most two tuple-indexing
compositions, the cost of \texttt{power(n)} is $O(\log n)$ compositions of
$O(|X|)$ work each --- so the factorial powers of the next paragraph are
computable even where $|X|!$ is astronomically large as an integer, since
only its bit pattern is traversed.

\texttt{factorial\_idempotent\_power()} returns
\texttt{power(factorial(len(carrier)))} --- the executable form of the safe
finite power $f^{|X|!}$ of Chapter~5. \texttt{is\_idempotent()} tests
\texttt{compose(arrow, arrow) == arrow}, an extensional comparison in which
the composite's fresh presentation name is irrelevant by
Invariant~\ref{ch09:inv:extensional}.

\begin{remark}[The docstring as a claim boundary]\label{ch09:rem:power-docstring}
The \texttt{power} docstring carries a sentence that belongs to the claim
discipline as much as to the API documentation: ``This operation is
deliberately absent from \texttt{SemigroupoidPath}. It does not license
powers of heterogeneous or time-varying paths.'' The prohibition that
Chapter~6 develops mathematically is thus stated at the exact point in the
code where a maintainer contemplating ``adding \texttt{power} to paths for
convenience'' would read it.
\end{remark}

\subsection{\texttt{SemigroupoidPath}: heterogeneous composites without powers}\label{ch09:sec:path}

\begin{construction}[The path type]\label{ch09:con:path}
\texttt{SemigroupoidPath} is a frozen, slotted dataclass wrapping
\texttt{arrows: tuple[FiniteArrow, ...]}. Validation enforces, in order:
the sequence passes the ordered-input guard; it is nonempty (``a
semigroupoid path is nonempty''); every member is a \texttt{FiniteArrow};
and every adjacent pair matches, \texttt{before.target == after.source}
under full object equality, with violations raising
\texttt{CompositionTypeError} whose message prints both segments' type
signatures under the prefix ``ill-typed path segment.''
\end{construction}

The surface is: a classmethod \texttt{from\_iterable(arrows)} that simply
delegates to the constructor (its docstring --- ``Build from an explicitly
ordered, replayable sequence'' --- restates the input contract);
\texttt{source} and \texttt{target} properties projecting the first arrow's
source and the last arrow's target; \texttt{composite()}, which folds
\texttt{compose} left to right over the arrows, naming intermediate results
\texttt{path-composite[k]} for the step count $k$, and returns a final
arrow renamed \texttt{path-composite[n]} for the path length $n$; and
\texttt{append(arrow)}, which returns a \emph{new} path with the arrow
adjoined (revalidating the junction through the constructor --- there is no
in-place mutation anywhere in the package).

\begin{invariant}[No power operation on paths]\label{ch09:inv:no-path-power}
\texttt{SemigroupoidPath} defines no \texttt{power} method, no
\texttt{\_\_pow\_\_}, no \texttt{\_\_mul\_\_}, and no conversion to
\texttt{Endomap}. A path whose endpoints happen to coincide still cannot be
raised to a power without the caller explicitly extracting
\texttt{composite()} and wrapping it in \texttt{Endomap} --- at which point
the endomap validator independently certifies that the composite is endo.
The invalid inference from ``sequence of generators'' to ``element of one
finite monoid'' is thereby not merely discouraged but unwritable in the
type vocabulary the module exports, exactly as PO-TKS8-CORE-03 requires.
\end{invariant}

\subsection{Design decisions visible in \texttt{category.py}}\label{ch09:sec:cat-design}

Gathering the threads, six decisions are legible in the module's source and
worth recording as decisions rather than accidents:

\begin{enumerate}
\item \emph{Frozen, slotted dataclasses everywhere.} Immutability makes
      every value hashable and safely shareable; slots foreclose attribute
      injection; and both together mean an object or arrow, once
      validated, can never drift from its validated state.
\item \emph{Validation in \texttt{\_\_post\_init\_\_}, never later.} All
      preconditions are checked at construction; no method needs to
      re-validate its own receiver. An invalid value is unrepresentable,
      not merely detected.
\item \emph{The equality asymmetry.} Object names are type-forming; arrow
      names are presentation-only. One \texttt{compare=False} keyword
      carries the whole extensionality discipline.
\item \emph{Two exceptions, mapped onto Python's own value/type split.}
      Malformed data is a \texttt{ValueError} subclass; ill-typed use of
      well-formed data is a \texttt{TypeError} subclass.
\item \emph{The \texttt{bool} guard.} Both the arrow mapping validator and
      the exponent validator explicitly reject \texttt{bool} before
      accepting \texttt{int}, closing the subclass loophole through which
      \texttt{True} would otherwise be read as the index $1$.
\item \emph{Bounded names.} Presentation-name construction is $O(\log n)$
      in any exponent via the compact label helper, so no name-building
      code path can dominate, overflow, or hit interpreter conversion
      limits --- and since names are outside equality, the bounding is
      free of semantic cost.
\end{enumerate}

\section{\texttt{deterministic\_lift.py}: the standard-basis lift}\label{ch09:sec:lift}

\subsection{Purpose and representation policy}

\texttt{deterministic\_lift.py} realizes the lift functor of Chapter~7: for
a typed arrow $f : X \to Y$, the linear map $A_f$ on standard bases with
$A_f e_x = e_{f(x)}$. The module docstring states both the norm fact ---
``the exact operator norm is $\sqrt{\max_y |f^{-1}(y)|}$ (zero for an empty
domain)'' --- and the representation policy: ``The implementation exposes
exact integer-valued squared data wherever possible; only
\texttt{operator\_norm} converts to floating point.''

That policy is the module's central design decision. Every matrix the
module returns --- the lift matrix itself, both Gram matrices, the support
projector --- is a tuple of tuples of Python \texttt{int}s. Every scalar
invariant except one --- squared norm, maximum fiber size, trace, rank ---
is an exact integer. The single floating-point value in the module is the
property \texttt{operator\_norm}, computed as \texttt{sqrt} of the exact
integer \texttt{operator\_norm\_squared}; a caller who needs exactness
compares squared norms and never calls it. This is what makes the
$1{,}678$-pair functoriality check of Chapter~10 an \emph{exact} matrix
comparison with no tolerance parameter: there is no rounding anywhere to
tolerate.

\subsection{Construction and the applied maps}

\begin{construction}[The lift wrapper]\label{ch09:con:lift}
\texttt{DeterministicLift} is a frozen, slotted dataclass wrapping
\texttt{arrow: FiniteArrow}; its validator rejects non-\texttt{FiniteArrow}
values with \texttt{ValidationError}. There is no other precondition ---
every typed arrow lifts, including arrows with empty source or empty
target.
\end{construction}

\texttt{matrix()} returns the $|Y| \times |X|$ zero-one matrix as a tuple
of row tuples: entry $(r, c)$ is \texttt{int(mapping[c] == r)}. Each column
$c$ therefore contains exactly one $1$, in row \texttt{mapping[c]} --- the
column-selection form derived in Chapter~7.

\texttt{apply(vector)} applies $A_f$ to source coordinates. Its guard
repeats the ordered-sequence check inline (rejecting strings, bytes,
mappings, sets, and non-sequences with ``source coordinates must be an
ordered sequence''), then checks the length against $|X|$ with a message
reporting both expected and received counts. The computation is a
scatter-add:

\begin{quote}
\texttt{output[mapping[source\_index]] += value}
\end{quote}

over an output list initialized to zeros of length $|Y|$. Coordinates on
colliding inputs accumulate --- this is precisely the summation
$ (A_f v)_y = \sum_{x \in f^{-1}(y)} v_x$, and it is why a constant map's
lift is noncontractive: mass concentrates rather than spreads.

\texttt{adjoint\_apply(vector)} applies the conjugate transpose $A_f^{*}$
to target coordinates, with the mirrored guard and length check against
$|Y|$. The computation is a gather:

\begin{quote}
\texttt{tuple(vector[target\_index] for target\_index in mapping)}
\end{quote}

--- entry $x$ of the result is coordinate $f(x)$ of the input, which is
$(A_f^* w)_x = w_{f(x)}$. Because all matrix entries are real integers,
conjugation is trivially the transpose; the docstring nonetheless says
``conjugate transpose,'' keeping the Hilbert-space reading in view. Note
that the pair \texttt{apply}/\texttt{adjoint\_apply} lets a test verify
adjointness numerically --- $\langle A_f v, w\rangle = \langle v, A_f^* w
\rangle$ --- without ever materializing a matrix.

\subsection{Fibers and the derived operators}

The property \texttt{fiber\_sizes} computes, in one pass over the mapping,
the tuple of fiber cardinalities $\bigl(|f^{-1}(y)|\bigr)_{y \in Y}$: a
zero-initialized list of length $|Y|$, incremented at each image index.
Everything else in the module is a pure function of this tuple and the
mapping:

\begin{itemize}
\item \texttt{max\_fiber\_size} is \texttt{max(fiber\_sizes, default=0)};
      the \texttt{default} argument realizes the empty-target convention
      --- and since an empty \emph{source} yields an all-zero tuple, the
      empty-domain value is also $0$, matching the docstring's
      parenthesis.
\item \texttt{operator\_norm\_squared} returns \texttt{max\_fiber\_size}
      and is documented as ``the exact value
      $\lVert A_f\rVert^2 = \max_y |f^{-1}(y)|$'' --- the norm theorem of
      Chapter~7 hard-coded as a property. Note the implementation does not
      \emph{compute} a spectral radius; it computes the right-hand side of
      a theorem proved off-line. The structural tests of Chapter~10 are
      what tie this shortcut to matrix reality on the bounded family.
\item \texttt{operator\_norm} is the lone float:
      \texttt{sqrt(operator\_norm\_squared)}.
\item \texttt{input\_gram\_matrix()} returns $A_f^* A_f$ directly from the
      collision structure: entry $(x, x')$ is
      \texttt{int(mapping[x] == mapping[x'])}. No matrix multiplication is
      performed --- the code implements the \emph{lemma} that the Gram
      entry indicates fiber-mates, not the product definition.
\item \texttt{output\_gram\_matrix()} returns $A_f A_f^*$ as the diagonal
      matrix of fiber sizes.
\item \texttt{output\_support\_projector()} returns the diagonal zero-one
      matrix marking nonempty fibers --- ``the orthogonal projector onto
      $\im(A_f)$ in the codomain.''
\end{itemize}

\begin{remark}[Three operators kept distinct]\label{ch09:rem:three-ops}
Chapter~7 warned that informal usage conflates $A_f^* A_f$, $A_f A_f^*$,
and the support projector. The module's answer is architectural: they are
three differently named methods with three different docstrings, and the
second and third agree only when every fiber has size at most one. A caller
cannot obtain one while asking for another, and the test suite pins each
method to its own fixture.
\end{remark}

\subsection{Scalar invariants and predicates}

The remaining properties read off the implemented consequences of the norm
theory:

\begin{center}
\begin{tabular}{@{}lll@{}}
\toprule
Property & Value & Justification in code \\
\midrule
\texttt{trace\_input\_gram} & $|X|$ & \texttt{len(arrow.source)} \\
\texttt{rank} & $|\im f|$ & count of nonzero fiber sizes \\
\texttt{is\_injective} & $f$ injective & all fiber sizes $\le 1$ \\
\texttt{is\_isometry} & $A_f^*A_f = I$ & delegates to \texttt{is\_injective} \\
\texttt{is\_partial\_isometry} & --- & delegates to \texttt{is\_injective} \\
\texttt{is\_idempotent} & $A_f^2 = A_f$ & see below \\
\texttt{is\_orthogonal\_projector} & $A_f = A_f^* = A_f^2$ & endomap and identity mapping \\
\bottomrule
\end{tabular}
\end{center}

\texttt{trace\_input\_gram} does not sum the Gram diagonal; it returns the
source size directly, implementing the identity
$\trace(A_f^* A_f) = |X|$ rather than computing its left side ---
the same theorem-as-shortcut pattern as the norm, and tested the same way.
\texttt{rank} counts nonempty fibers, implementing
$\rk(A_f) = |\im f|$. \texttt{is\_isometry}'s docstring records the
equivalence it relies on (``whether $A_f^* A_f$ is the identity
(equivalently, $f$ injective)''), and \texttt{is\_partial\_isometry}'s
docstring notes the special situation of total deterministic lifts: ``for a
total standard lift, equivalently whether $f$ is injective.'' For general
operators partial isometry is strictly weaker than isometry; for
column-selection matrices the two coincide, and the code documents that it
is using the coincidence rather than the general definition.

\texttt{is\_idempotent} first returns \texttt{False} for non-endomaps ---
$A_f^2$ is not even typeable on one space otherwise --- and then checks

\begin{quote}
\texttt{all(mapping[image] == image for image in mapping)}
\end{quote}

which says exactly that $f$ fixes every point of its image, the
combinatorial characterization of $f \circ f = f$. No matrix is squared.

\texttt{is\_orthogonal\_projector} is the most carefully documented
predicate in the module. Its docstring states the fact the implementation
encodes: ``A total deterministic lift on one declared carrier is an
orthogonal projector only for the identity map. Non-endomaps are not
operators on one Hilbert space and therefore return \texttt{False}.'' The
check is accordingly \texttt{is\_endomap} together with the mapping equaling
\texttt{tuple(range(len(source)))}. The mathematics behind the docstring is
Chapter~7's: an orthogonal projector must be self-adjoint, a zero-one
column-selection matrix is symmetric only when it is a permutation matrix
fixed by transposition on its support, and combining self-adjointness with
idempotence forces the identity. The constant-map fixture --- idempotent,
norm $\sqrt{2}$, not self-adjoint --- is exactly the witness that
\texttt{is\_idempotent} and \texttt{is\_orthogonal\_projector} must not be
merged, and the permanent regression test holds them apart.

\subsection{Error taxonomy of the lift module}

The module introduces no new exception types. All rejections raise
\texttt{ValidationError} imported from \texttt{category.py}: wrapping a
non-arrow, passing an unordered container as a coordinate vector, or
passing a vector of the wrong length. There is no
\texttt{CompositionTypeError} here because the module performs no
composition --- functoriality is a property the \emph{tests} check by
composing arrows in \texttt{category.py} and comparing matrices, not an
operation the lift class itself exposes. The lift module consequently has
no way to combine two lifts incorrectly: it has no way to combine two
lifts at all.

\section{\texttt{normal\_forms.py}: the shortlex closure engine}\label{ch09:sec:nf}

\subsection{Purpose}

\texttt{normal\_forms.py} implements the denotational normal-form catalog
of Chapter~8: given a finite, explicitly ordered set of typed generators
over a finite object set, compute for every reachable finite-function
denotation the shortlex-least well-typed word denoting it. The module
docstring compresses the claim discipline into three sentences: no rewrite
system is assumed; finiteness of the carriers makes the denotation closure
finite; and ``construction either completes or raises a resource guard
exception, so an incomplete table is never presented as canonical.''

\subsection{Error taxonomy}

Two exception types are defined here, again with meaningful base classes:

\begin{center}
\begin{tabular}{@{}lll@{}}
\toprule
Exception & Base class & Meaning \\
\midrule
\texttt{ClosureLimitExceeded} & \texttt{RuntimeError} & a resource guard tripped; no partial catalog \\
\texttt{UnreachableDenotationError} & \texttt{LookupError} & arrow not generated by the declared generators \\
\bottomrule
\end{tabular}
\end{center}

The docstring of \texttt{ClosureLimitExceeded} is a policy statement:
``Raised instead of returning a partial normal-form catalog.''
\texttt{UnreachableDenotationError} subclasses \texttt{LookupError} because
it is semantically a failed lookup --- the caller asked for the
representative of a denotation the closure does not contain --- and
distinguishing it from \texttt{ValidationError} lets callers separate
``you asked a malformed question'' from ``you asked a well-formed question
whose answer is that the arrow is outside the generated subcategory.''
Malformed inputs to this module (blank labels, undeclared objects, unknown
generator labels, non-positive resource bounds) raise the imported
\texttt{ValidationError}, keeping the package-wide taxonomy at four
exception types total.

\subsection{\texttt{Generator} and \texttt{TypedWord}}

\texttt{Generator} is a frozen, slotted dataclass pairing a \texttt{label}
(a \texttt{str} whose \texttt{strip()} must be nonempty) with an
\texttt{arrow} (a \texttt{FiniteArrow}). The label is the unit of word
spelling; the arrow is its denotation.

\texttt{TypedWord} is a frozen, slotted dataclass with fields
\texttt{start: FiniteObject} and \texttt{labels: tuple[str, ...]}
(defaulting to the empty tuple). Its validator checks that \texttt{start}
is a \texttt{FiniteObject}, passes the labels through the ordered-input
guard, and requires every label to be a nonempty string. The class
docstring records a division of labor that the constructor cannot see
past: ``The constructor validates syntax. Categorical well-typedness is
relative to a generator catalog and is certified by
\texttt{NormalFormCatalog.evaluate}.'' A \texttt{TypedWord} standing alone
is a spelling; only a catalog can say whether its labels resolve and
compose. Carrying the start object explicitly is what makes the empty word
meaningful --- the empty word at $X$ denotes $\id_X$, and without a
declared start there would be one empty word instead of one per object.

\subsection{\texttt{NormalFormCatalog}: constructor validation}

The catalog constructor takes an ordered sequence of objects, an ordered
sequence of generators, and two keyword-only resource bounds:
\texttt{max\_arrows} (default $100{,}000$) and \texttt{max\_frontier}
(default \texttt{None}, in which case it is set equal to
\texttt{max\_arrows}). Validation is extensive and worth enumerating
completely, because each rule closes a distinct determinism or
well-formedness hole:

\begin{enumerate}
\item objects and generators each pass the ordered-input guard --- their
      declared order is data (it fixes shortlex ranks), so unordered
      containers are rejected at the door;
\item the object sequence is nonempty (``a normal-form catalog needs at
      least one object'');
\item every declared object is a \texttt{FiniteObject}; every declared
      generator is a \texttt{Generator};
\item the objects are pairwise distinct \emph{as objects}, and separately
      their \emph{names} are pairwise distinct --- the second check is
      strictly stronger and exists so that error messages and external
      references can use names unambiguously;
\item \texttt{max\_arrows} and \texttt{max\_frontier} are positive
      \texttt{int}s, with the \texttt{bool} subclass explicitly excluded,
      as in the arrow and exponent validators;
\item generator labels are pairwise distinct --- a duplicated label would
      make word spellings ambiguous;
\item every generator's source and target appear among the declared
      objects (``generator ... references an undeclared object'') --- the
      catalog's object list is a closed world, and the closure can never
      escape it.
\end{enumerate}

After validation the constructor builds four private tables --- label to
generator, label to rank (position in the declared generator order),
object to rank (position in the declared object order), and the two
closure dictionaries mapping denotation keys to representative words and
to arrows --- and then immediately runs \texttt{\_build}. A catalog is
therefore never observable in a half-constructed state: by the time the
constructor returns, the closure is complete, or an exception has
propagated and there is no catalog at all. This constructor-runs-closure
design is the object-level form of the fail-closed policy.

\subsection{Shortlex keys and evaluation}

\texttt{shortlex\_key(word)} maps a word to the triple

\[
(\;|w|,\; \text{rank sequence of } w,\; \text{rank of start object}\;)
\]

realized as \texttt{(len(labels), ranks, start\_rank)}. Unknown labels and
undeclared start objects raise \texttt{ValidationError} with the offending
name in the message. Python's lexicographic tuple comparison then
implements exactly the shortlex order of Chapter~8 --- length first, then
generator ranks positionwise --- extended by a final tie-break on the
start-object rank, which orders the empty words of distinct objects and
never separates two words of the same start and spelling. This start-rank
component is a totalizing convenience for the catalog-wide sorted listing
(\texttt{normal\_forms} below), not a change to the per-denotation
minimization: all words denoting one fixed arrow share the arrow's source,
hence share the tie-break.

\texttt{evaluate(word)} certifies well-typedness by construction: it
starts from \texttt{identity(word.start)} (rejecting undeclared starts),
then left-folds \texttt{compose(generator.arrow, result)} over the labels,
resolving each label through the catalog table (unknown labels raise
\texttt{ValidationError}) and letting \texttt{compose} itself raise
\texttt{CompositionTypeError} at the first ill-typed junction. The result
is repackaged under the presentation name
\texttt{word-evaluation[n]}. Evaluation is thus the one gate through which
a word acquires a denotation, and it inherits every type check of
Section~\ref{ch09:sec:compose} rather than re-implementing any.

Three query methods complete the public surface.
\texttt{normalize(word)} evaluates the word and returns the stored
representative for its denotation key --- by completeness of the closure,
any word that evaluates successfully over declared objects and generators
has a representative, so \texttt{normalize} on well-formed input does not
miss. \texttt{representative\_for(arrow)} looks up an arbitrary arrow's
denotation key directly and raises \texttt{UnreachableDenotationError}
(naming the arrow) when the arrow is outside the generated closure ---
this is the method a downstream consumer uses to ask ``is this function
expressible in the declared generators, and if so, what is its canonical
spelling?'' The property \texttt{denotation\_count} reports the closure
size, and \texttt{normal\_forms} returns all representatives sorted by
\texttt{shortlex\_key} --- a deterministic, replayable listing of the
entire generated subcategory.

\subsection{The closure engine \texttt{\_build}}\label{ch09:sec:build}

The private method \texttt{\_build} is the algorithmic heart of the module
--- the executable form of the representative-first closure of Chapter~8.
Its machinery:

\begin{itemize}
\item a min-heap \texttt{frontier} of quadruples
      \texttt{(shortlex\_key, serial, word, arrow)};
\item a monotone counter \texttt{serial} from \texttt{itertools.count},
      whose value is the second heap component --- it breaks exact key
      ties by insertion order and, decisively, prevents the heap from ever
      comparing the third and fourth components, since \texttt{TypedWord}
      and \texttt{FiniteArrow} admit equality but not ordering;
\item a set \texttt{queued\_denotations} of denotation keys currently in
      the frontier, used to deduplicate candidates before they are pushed.
\end{itemize}

The run has two phases. \emph{Seeding:} for each declared object, in
declared order, the empty word and identity arrow are pushed --- each push
preceded by a frontier-capacity check that raises
\texttt{ClosureLimitExceeded} if the queued-denotation count has reached
\texttt{max\_frontier}. Identities are thus in the closure by fiat, with
the empty word as their (trivially shortlex-least) spelling.
\emph{Exploration:} while the frontier is nonempty, pop the least
quadruple; discard its key from the queued set; skip it if the key is
already represented (a defensive guard --- the push-side deduplication
already prevents duplicate enqueueing, and the pop-side check makes the
invariant locally evident); check the \texttt{max\_arrows} guard
\emph{before} inserting, raising \texttt{ClosureLimitExceeded} if the
accepted count has reached the cap; record the word as the representative
and the arrow as the stored denotation; then attempt to extend by every
generator in declared order, skipping generators whose source does not
equal the current arrow's target, building the extended word (labels plus
the generator's label) and the extended arrow (via \texttt{compose}, under
the internal name \texttt{normal-form-closure-candidate}), and pushing the
extension --- with the same frontier-capacity check --- only if its
denotation key is in neither the representative table nor the queued set.

\begin{invariant}[First insertion is shortlex-least]\label{ch09:inv:first-least}
Words are popped in nondecreasing shortlex order (heap order on the key,
with serial tie-break), and every word is pushed with the key
\texttt{shortlex\_key} assigns it. By the pruning-soundness argument of
Chapter~8, the first word popped for a given denotation key is the
shortlex-least well-typed word with that denotation, and pruning later
words with the same key removes no denotation and no least
representative. The dictionary insertion in \texttt{\_build} is therefore
the whole of the minimization: no post-processing compares or replaces
representatives.
\end{invariant}

\begin{construction}[Resource guards]\label{ch09:con:guards}
Two independent caps bound the run. \texttt{max\_arrows} bounds the number
of accepted denotations: the guard fires on the pop side, before the
insertion that would exceed the cap. \texttt{max\_frontier} bounds the
number of distinct denotations simultaneously queued: the guard fires on
the push side, before the push that would exceed the cap --- including
during seeding, so a catalog declaring more objects than
\texttt{max\_frontier} fails immediately rather than after partial work.
Both guards raise \texttt{ClosureLimitExceeded} with messages that name
the tripped bound and repeat the policy: ``no partial normal-form catalog
was returned.'' Because the guards raise out of the constructor, the
partially built dictionaries are unreachable --- there is no API through
which a caller could observe the truncated state.
\end{construction}

\begin{remark}[Why deduplicate the frontier by denotation]\label{ch09:rem:dedupe}
The push-side check --- skip extensions whose key is already represented
\emph{or already queued} --- is a completeness-preserving optimization
justified by the same pruning argument: if some word for a denotation is
already queued, the queued one has a shortlex key at most that of the
candidate arriving later (candidates are generated in nondecreasing prefix
order, and equal keys are resolved by the earlier serial), so discarding
the newcomer never discards a least representative. Operationally it also
makes \texttt{max\_frontier} meaningful as a bound on \emph{distinct}
pending denotations rather than on raw queue length, tightening the
resource semantics the guard enforces.
\end{remark}

\subsection{Design decisions visible in \texttt{normal\_forms.py}}

Four decisions deserve explicit record. First, \emph{completion in the
constructor}: there is no lazy or incremental mode, so every reachable
catalog object is a completed catalog, and the fail-closed policy needs no
cooperation from callers. Second, \emph{denotation keys as the sole
identity}: the closure tables are keyed by
\texttt{arrow.denotation\_key} --- the extensional triple --- so all the
presentation names generated during search
(\texttt{normal-form-closure-candidate}, \texttt{word-evaluation[n]}) are
invisible to the mathematics, as they must be. Third, \emph{ranks from
declared order only}: generator rank and object rank are positions in the
caller's sequences; nothing is sorted by name, hash, or carrier content,
so the canonical forms are exactly as replayable as the caller's declared
order --- which Chapter~8's claim boundary makes versioned data. Fourth,
\emph{the serial tie-break}: a two-line idiom (\texttt{count()} plus a
tuple slot) that simultaneously guarantees heap totality without ordering
instances and makes exploration order fully deterministic given the
declared orders --- determinism again pushed into structure rather than
left to accident.

\section{Manifest binding of the implementation}\label{ch09:sec:manifest}

\subsection{The artifact manifest}

The prototype is bound by the artifact manifest
\texttt{TKS\_V8\_CORE\_PROTOTYPE\_MANIFEST\_v0.1.json} (artifact id
\textsc{tks-v8-typed-core-prototype}, artifact version 0.1.0, schema
version 0.1). The manifest records the artifact's status block ---
\textsc{gate 1 prototype / in review / not frozen}, admission ring 4,
license status \textsc{not licensed} --- its provenance (the governing
mathematical red-team report, the theorem ledger, and the three theorem
record ids \textsc{thm-lift-norm-fiber-formula},
\textsc{thm-lift-occupation-counterexample}, and
\textsc{thm-finite-idempotent-power}), its integration boundary, its
formal-verification and validation snapshots (documented in Chapters~10
and~11), its claim records, and --- the subject of this section --- its
\emph{artifact bindings}: SHA-256 digests and byte counts for every file
the artifact comprises.

\subsection{The eight code bindings}

Four implementation files and four test modules are hash-bound with the
following roles, quoted from the manifest:

\begin{center}
\small
\begin{tabular}{@{}lp{5.4cm}r@{}}
\toprule
Path & Role & Bytes \\
\midrule
\texttt{src/tks\_v8\_core/\_\_init\_\_.py} & package exports & 1{,}018 \\
\texttt{src/tks\_v8\_core/category.py} & typed finite-function category & 9{,}850 \\
\texttt{src/tks\_v8\_core/deterministic\_lift.py} & deterministic standard-basis lift & 5{,}483 \\
\texttt{src/tks\_v8\_core/normal\_forms.py} & denotational shortlex closure & 9{,}889 \\
\texttt{tests/tks\_v8\_core/test\_category.py} & category and power regressions & 5{,}588 \\
\texttt{tests/tks\_v8\_core/test\_deterministic\_lift.py} & lift structure and functor-law regressions & 9{,}142 \\
\texttt{tests/tks\_v8\_core/test\_manifest\_boundaries.py} & claim and integration boundary regression & 1{,}965 \\
\texttt{tests/tks\_v8\_core/test\_normal\_forms.py} & normal-form correctness and resource regressions & 5{,}672 \\
\bottomrule
\end{tabular}
\end{center}

The SHA-256 digests, exactly as recorded:

\begin{itemize}
\item \texttt{\_\_init\_\_.py}: \\
      \texttt{6D7559758D5CB9A2CE4A9CD0EC152D0C8A7B9057708429571FEC1374639CDB77}
\item \texttt{category.py}: \\
      \texttt{B69F494814E9D9CD54C80ADD8D6AF850C5A8F762B1A0C6DFC85435BA4458D947}
\item \texttt{deterministic\_lift.py}: \\
      \texttt{A4649F62B44ED0DE820316E930958D2E32E143406DC1F9BC169A8EC533095115}
\item \texttt{normal\_forms.py}: \\
      \texttt{7C85D6BA77E39F953F59916C907E4D7C9C0564C7AB0E89CAE05087188BC596E5}
\item \texttt{test\_category.py}: \\
      \texttt{0D8592E1231FA02263E63496BE0774441610D3E11A17A6D18DF39DBBCBD789C2}
\item \texttt{test\_deterministic\_lift.py}: \\
      \texttt{D14126D80A056ED222DB3DFC58BEA4DFFF0E6466D31C43BC8D7B2F78CE509D01}
\item \texttt{test\_manifest\_boundaries.py}: \\
      \texttt{9D3D1DEA155A5E1F8A13638B915EEF28C4DF51A90226B5826BA862BF10F2E097}
\item \texttt{test\_normal\_forms.py}: \\
      \texttt{A4DE6902CC4AD986AAAB1A448473C3D72AE528907577E6F6F1CE328E0C396F9B}
\end{itemize}

Beyond the eight code files, the manifest binds four further artifacts:
the formal specification document
\texttt{TKS\_V8\_CORE\_FORMAL\_SPEC\_AND\_PROOF\_OBLIGATIONS\_v0.1.md}
(role: formal specification and explicit obligations; SHA-256
\texttt{F740BA05AB7ADB2DCDF55E271016D127E45073B015497DE8982D44A49F0EA2C4},
10{,}917 bytes), the integrated Lean receipt
\texttt{LEAN\_PROOF\_MANIFEST\_v0.1.json}
(\texttt{7C27E50DCD6E45177F279A3B57BCF9B5C27B9D9FE878D34A4AE66742D332EC4A},
5{,}959 bytes), the lift-functor formal subreceipt
\texttt{LEAN\_DETERMINISTIC\_LIFT\_FUNCTOR\_MANIFEST\_v0.1.json}
(\texttt{BEED66D60666D48EF33204779D8DCE0557AD43245E5D89CD63691CC97492C3A2},
4{,}509 bytes), and the verifier script
\texttt{scripts/verify\_tks\_v8\_gate1.py} (role: hash and executable test
verifier; SHA-256
\texttt{E4FF72C5980059BE9359509827CE4D96F39D12765A003CD25AD5063A0590BC28},
4{,}230 bytes). The specification of Chapters~1--8, the Lean receipts of
Chapter~11, and the verifier of Chapter~10 are thus all inside the same
binding envelope as the code they govern: any byte change to any of the
twelve files is detectable against the manifest.

\subsection{What the binding does and does not establish}

\begin{protocol}[Hash-binding discipline]\label{ch09:proto:binding}
The binding envelope supports one operation, implemented by
\texttt{scripts/verify\_tks\_v8\_gate1.py} and detailed in Chapter~10:
recompute the SHA-256 of every bound path, compare against the manifest,
and fail on any mismatch or missing file before rerunning the executable
suite. A successful run establishes that the bytes examined are the bytes
the manifest describes and that those bytes pass their own tests ---
nothing further.
\end{protocol}

\begin{boundary}[The manifest does not hash itself]\label{ch09:bdy:self-hash}
The manifest's own \texttt{binding\_boundary} field states the limit of the
scheme in its own words: ``The manifest does not hash itself. A later
reviewed commit or acceptance receipt must bind this manifest SHA-256.
Until then this is an uncommitted Gate~1 candidate, not a freeze
receipt.'' The point is structural, not procedural: a document cannot
contain its own digest, so the chain of custody necessarily terminates
outside the manifest --- in a reviewed commit, a signed receipt, or some
other artifact that records the manifest's hash under independent
authority. Until such an outer binding exists, an adversary (or an
accident) that rewrites the manifest \emph{and} the files consistently
would pass verification; the manifest therefore certifies internal
consistency of the snapshot, and the snapshot's authority awaits the
outer commitment. Recording this boundary here, in the manifest's own
language, is part of documenting the artifact as it is.
\end{boundary}

\section{Summary}\label{ch09:sec:summary}

The implementation is small by design: three modules, roughly twenty-five
kilobytes, four exception types, no dependencies beyond the standard
library. Its character is uniform across the modules. Validation is total
and constructor-resident, so invalid values are unrepresentable rather than
detected. Identity discipline is carried by ordinary equality on frozen
data, with exactly one asymmetry --- objects own their names, arrows do not
--- doing the work of the extensionality theory. Determinism is enforced at
the input boundary by a single ordered-sequence guard used everywhere order
is semantic. Theorems proved in the mathematical chapters appear in the
code as directly implemented right-hand sides --- the norm as a maximum
fiber size, the trace as a source cardinality, idempotence as an
image-fixing check --- with the bounded-exhaustive suites of Chapter~10
tying the shortcuts back to matrix reality. Prohibitions appear as absent
methods (\texttt{power} on paths) and as predicates that refuse to conflate
(\texttt{is\_idempotent} against \texttt{is\_orthogonal\_projector}).
Resource limits fail closed, out of constructors, leaving no observable
partial state. And the whole --- sources, tests, specification, Lean
receipts, verifier --- sits inside one hash-binding envelope whose sole
acknowledged gap, self-binding, is documented in the artifact's own text
and assigned to a later reviewed commit. Chapter~10 examines how this code
is tested; Chapter~11, what a proof assistant has and has not checked about
its mathematics; Chapter~12, what may and may not yet be built on it.
